% !TEX root = ../double_trouble_supp.tex

\subsection{Sampling from the proposed two population model via (truncated) Poisson point processes}
\label{app:sample_from_pois}

We here provide a practical algorithm to generate synthetic data from the proposed hierarchical model introduced in Equation (3.1) where the L\'{e}vy rate measure is given in \cref{eq:proposed_levy}.

At a high level, there exist two main strategies to generate data from nonparametric hierarchical models like the one in Equation (3.1). The first is to devise an exact ``marginal'' representation, where the infinitude of latent atoms is integrated out and observations are generated by leveraging an iterative predictive ``marginal'' scheme (e.g., the Chinese restaurant process, the Indian buffet process and related ``urn-schemes''). The other approach is to rely on an approximate ``conditional'' representation, in which the infinitude of latent atoms is approximated via a truncated representation of the infinite measure \citep{campbell2019truncated}, and observations are generated conditionally i.i.d.\ on this finite representation.

In principle we can derive a marginal representation of the proposed process following \citet{Broderick2018}. However such marginal representation would require us to solve many numerical integrals for every sample to determine whether a pre existing variant should be present in a follow up sample. 
To avoid this computational overhead, we decide to adopt the approximate (truncated) approach: we draw a truncated, approximate Poisson point process and obtain a sample $X_{\popidx, \unitidx}$ by repeatedly sampling Bernoulli trials using the sampled variant frequencies. 


A general approach to sample an inhomogeneous Poisson point process with an arbitrary target density is to find a ``proposal'' rate measure dominating the target that is easy to sample from (e.g., a homogeneous one), and then accept a point with probability equal to the ratio between the target density and the proposal density at that point \citep{saltzman2012simulating}. 
For our specific proposal, we lower bound the domain of our rate measure, enforcing variant frequencies to be larger than $10^{-10}$. This approximation is not concerning, since the probability of sampling a variant with a frequency smaller than this value in 500 samples as in our simulation setting is practically equal to $0$.

To sample variants frequencies from our (truncated) rate measure, we choose a piecewise constant proposal rate measure that dominates our (truncated) rate. To do so, we span a (log-scaled) mesh using $\{10^{-10},10^{-9},\dots, 10^{0}\} \times \{10^{-10},10^{-9},\dots, 10^{0}\}$. These grids are disjoint so we can sample points within each grid independently due to the properties of PPP. In each cell $\mathcal{G}$ of the grid, we sample a homogeneous Poisson point process with rate equal to the maximum value of our rate measure within in that grid cell ($\lambda_{m,\mathcal{G}}=\max_{\vecvariantfreq{}\in \mathcal{G}} \ratemeasproposed(\vecvariantfreq{})$). To do so, we first sample number of points by sampling a Poisson random variable whose mean is the size of the grid times the rate $L_{\mathcal{G}}\sim \Pois(\lambda_{m,\mathcal{G}}\cdot|\mathcal{G}|)$, then for each point, we sample its location uniformly in $\mathcal{G}$. Denote these samples as $\{\vecvariantfreq{\mathcal{G},\variantidx}\}_{\variantidx=1}^{L_{\mathcal{G}}}$. To get a sample from the truncated proposed rate measure within that grid $\mathcal{G}$, and we accept a sample at $\vecvariantfreq{\mathcal{G},\variantidx}$ with probability $\ratemeasproposed(\vecvariantfreq{\mathcal{G},\variantidx})/\lambda_{m,\mathcal{G}}$. We sample each grid $\mathcal{G}$ independently and combined all samples. Via this procedure, we sample an approximate realization our proposed PPP. Conditionally on these variants frequencies, we then sample the multi-population Bernoulli process $\mBeP$ of Equation (3.1). 




\subsection{d3BP}
To obtain a d3BP-Bernoulli process of size $\vecsamplesize = (\samplesize{1}, \samplesize{2})$ we first obtain $\samplesize{1} + \samplesize{2}$ samples from a 3BP-Bernoulli process using the Indian buffet representation \citep{Broderick2012,Masoero2021}. Then, we randomly split individuals into two groups of size $\samplesize{1}$ and $\samplesize{2}$. 

